% Part: lambda-calculus
% Chapter: introduction
% Section: syntax

\documentclass[../../../include/open-logic-section]{subfiles}

\begin{document}

\olfileid{lam}{int}{syn}
\olsection{لامبڈا کیلکولس دی نحو}

شروع وچ متغیراں $x$،$y$،$z$،~\dots دی اک لڑی تے کجھ مستقل
نشاناں $a$،$b$،$c$،~\dots لیندے آں۔ ٹرماں دا سیٹ استقرا نال ایویں متعین کردے آں:
\begin{enumerate}
\item ہر متغیر اک ٹرم اے۔
\item ہر مستقل اک ٹرم اے۔
\item جے $M$ تے $N$ ٹرم نیں،تاں $(MN)$ وی ٹرم اے۔
\item جے $M$ ٹرم اے تے $x$ متغیر اے،تاں $(\lambd[x][M])$ وی اک
  ٹرم اے۔
\end{enumerate}
$(MN)$ شکل والے ٹرماں نوں \emph{اطلاق}،تے
$(\lambd[x][M])$ شکل والے ٹرماں نوں \emph{تجرید} کہندے نیں۔

جیہڑے نظام وچ کوئی مستقل نشان نہ ہووے، اوہنوں \emph{خالص}
لامبڈا کیلکولس کہندے نیں۔اسی زیادہ تر خالص $\lambd$-کیلکولس
وچ کم کراں گے،اس لئی سارے چھوٹے لاطینی حروف متغیراں لئی ہون گے۔
وڈے لاطینی حروف ($M$،$N$،وغیرہ) $\lambd$-کیلکولس دے ٹرماں
لئی ورتاں گے۔

اسیں لکھت دے کجھ رواج ورتاں گے:

\begin{conv}
\begin{enumerate}
\item جدوں قوساں چھڈ دئیے، اطلاق کھبے توں سجے ول جُڑدا اے۔
  مثلاً جے $M$،$N$،$P$،تے $Q$ ٹرم نیں،تاں $MNPQ$،
  $(((MN)P)Q)$ دا مختصر روپ اے۔
\item قوساں چھڈن اُتے لامبڈا تجرید نوں ممکنہ سب توں وڈا دائرۂ
  اثر دتا جاندا اے۔مثلاً $\lambd[x][MNP]$ نوں
  $(\lambd[x][((MN)P)])$ پڑھیا جاندا اے۔
\item اک لامبڈا نال کئی متغیراں دی تجرید لکھی جا سکدی اے۔
  مثلاً $\lambd[xyz][M]$،
  $\lambd[x][\lambd[y][\lambd[z][M]]]$ دا مختصر روپ اے۔
\end{enumerate}
\end{conv}

مثلاً
\[
\lambd[xy][xxyx \lambd[z][xz]]
\]
ایس نال مراد اے:
\[
\lambd[x][\lambd[y][((((xx)y)x)(\lambd[z][(xz)]))]].
\]
ایہ رواج زبانی یاد کر لو۔ شروع وچ ایہ تُہانوں چکرا دین گے۔ عادت
پے جائے تاں ایہ قوساں دے ڈھیر نالوں گھٹ چکران گے۔

دو ٹرم جیہڑے صرف پابند متغیراں دے ناماں وچ فرق رکھدے نیں،اوہناں
نوں $\alpha$-ہم ارز کہندے نیں؛مثلاً $\lambd[x][x]$ تے
$\lambd[y][y]$۔ایہناں نوں ``اکّو'' ٹرم سمجھنا سہولت دی گل اے؛
دوجے لفظاں وچ،جدوں اسی کہندے آں کہ $M$ تے $N$ اکّو نیں،تاں
مراد اے کہ پابند متغیراں دے ناں بدل کے اوہ اکّو نیں۔ متغیر دی
جیہڑی آمد اپنے ناں والے $\lambd$ دے دائرۂ اثر وچ ہووے،اوہ
``پابند'' اے؛باقی ``آزاد'' اے۔پچھلی مثال وچ کوئی آزاد متغیر نہیں؛
پر
\[
(\lambd[z][yz])x
\]
وچ $y$ تے $x$ آزاد نیں،تے $z$ پابند اے۔

\end{document}
